\chapter{Implementazione}

\section{Classe template}

L'implementazione dell'algoritmo è stata effettuata in \texttt{C++}, attraverso l'utilizzo di una classe \texttt{template}, in modo da garantire l'utilizzo 
della classe con diversi tipi di chiavi e dati.
La classe è stata strutturata con degli enum in modo da andare a configurare opportunamente l'istanza con le strategie scelte per la hash table.

\begin{lstlisting}[language=C++, caption=HashTable class]
enum class OVERFLOW {
    OPEN_ADDRESSING_LINEAR_PROBING,
    OPEN_ADDRESSING_DOUBLE_HASHING,
    CHAINING,
    CUCKOO
};
enum class HASH {
    MODE_1,
    MODE_2
};
enum class RESIZE {
    DOUBLE,
    CONSTANT,
    DYNAMIC
};

template <typename Tkey, typename Tvalue>
class HashTable {
    public:
        HashTable(
            OVERFLOW overflow = OVERFLOW::OPEN_ADDRESSING_LINEAR_PROBING,
            HASH hash = HASH::MODE_1,
            RESIZE resize = RESIZE::DOUBLE,
            std::size_t table_length = 256,
            double max_load_factor = 0.75
        );
}
\end{lstlisting}


\section{Hash}

Per garantire un corretto funzionamento ed un corretto hashing al variare del tipo delle chiavi da inserire, queste vengono prima convertite 
attraverso l'apposita funzione \texttt{KeyConvert} e successivamente hashate e gestite secondo la strategia di collisione selezionata:

\begin{lstlisting}[language=C++, caption=KeyConvert function]
template<typename K>
struct KeyConvert {
    static std::size_t hash(const K& key) {
        if constexpr (std::is_integral_v<K> || std::is_same_v<K, char>) {
            return static_cast<std::size_t>(key);
        } else if constexpr (std::is_same_v<K, std::string>) {
            // stringhe
            std::size_t h = 0;
            for (char c : key) h = h * 31 + static_cast<unsigned char>(c);
            return h;
        } else {
            static_assert(sizeof(K) == 0, "Tipo non supportato");
        }
    }
};
\end{lstlisting}

Il progetto utilizza tre diverse hash function, ciascuna con uno scopo specifico:

\begin{enumerate}
    \item \textbf{Prima hash function (\texttt{h1})}:  
    Genera l'indice principale della tabella tramite l'operazione di \texttt{modulo} tra il valore hash della chiave e la lunghezza della tabella.

    \begin{lstlisting}[language=C++, caption=First hash function]
    template <typename Tkey, typename Tvalue>
    std::size_t HashTable<Tkey, Tvalue>::h1(const Tkey& key) const {
        std::size_t h = KeyConvert<Tkey>::hash(key);
        // modulo table_length
        return h % table_length;
    }
    \end{lstlisting}

    \item \textbf{Seconda hash function (\texttt{h2})}:  
    Può essere selezionata tramite apposito parametro durante la creazione della tabella.  
    Viene utilizzata sia per l'hashing di base alternativo sia come seconda hash function nella strategia Cuckoo.  
    Nel caso in cui sia selezionata per la seconda tabella della strategia Cuckoo, \texttt{h1} viene usata come fallback.
    La funzione moltiplica il valore hash della chiave per una costante molto grande 
    \texttt{a} e poi applica l’operazione di modulo con la lunghezza della tabella.  
    Questo approccio genera una distribuzione più uniforme degli indici rispetto a h1.

    \begin{lstlisting}[language=C++, caption=Second hash function]
    template <typename Tkey, typename Tvalue>
    std::size_t HashTable<Tkey, Tvalue>::h2(const Tkey& key) const {
        std::size_t h = KeyConvert<Tkey>::hash(key);
        constexpr std::size_t a = 11400714819323198485ull;
        return (h * a) % table_length;
    }
    \end{lstlisting}

    \item \textbf{Terza hash function (\texttt{h3})}:  
    Specifica per la strategia di double hashing, calcola l’incremento di probing in caso di collisioni.  
    La funzione utilizza un numero primo $R$ immediatamente inferiore alla dimensione della tabella per garantire che l’incremento e 
    la lunghezza della tabella siano coprimi, così da poter esplorare tutte le posizioni dell’array.

    \begin{lstlisting}[language=C++, caption=Hash function for double hashing]
    template <typename Tkey, typename Tvalue>
    std::size_t HashTable<Tkey, Tvalue>::h3(const Tkey& key) const {
        std::size_t h = KeyConvert<Tkey>::hash(key);
        std::size_t R = previousPrime(table_length - 1);
        return R - (h % R); 
    }
    \end{lstlisting}
\end{enumerate}

In questo modo, la tabella può supportare diversi tipi di chiave e diverse strategie di collisione, 
mantenendo le operazioni di inserimento, ricerca e cancellazione efficienti.


\section{Data Structures}

\subsection{Open Addressing}

Per l'open addressing, la struttura dati utilizzata per l'implementazione della hash table è un \texttt{std::vector<Slot>}, 
ovvero un vettore di oggetti \texttt{Slot}.  
Ogni \texttt{Slot} è caratterizzato da due campi per la chiave e il valore, più un campo \texttt{state} 
che indica se la cella è vuota, occupata o è stata rimossa.  
Questo campo aggiuntivo è fondamentale per preservare la logica di ricerca dell'open addressing anche dopo la cancellazione di elementi.

\begin{lstlisting}[language=C++, caption=Open Addressing Data Structure]
enum class SlotState { EMPTY, OCCUPIED, DELETED };

struct Slot {
    SlotState state = SlotState::EMPTY;
    Tkey key;
    Tvalue value;
};

std::vector<Slot> table_open;
\end{lstlisting}

---

\subsection{Chaining}

Per il chaining, la struttura dati utilizzata è un \texttt{std::vector<std::vector<std::pair<Tkey, Tvalue>>>},  
ovvero un vettore di vettori che rappresentano i bucket della tabella hash.  
Ogni bucket contiene una lista di coppie chiave-valore corrispondenti agli elementi presenti, permettendo di gestire più elementi nella stessa posizione senza perdere dati.

\begin{lstlisting}[language=C++, caption=Chaining Data Structure]
std::vector<std::vector<std::pair<Tkey, Tvalue>>> table_chaining;
\end{lstlisting}

---

\subsection{Cuckoo Hashing}

Per il Cuckoo Hashing, la struttura dati è costituita da due \texttt{std::vector<Slot>}, duplicando la struttura utilizzata per l'open addressing.  
Ogni \texttt{Slot} contiene la chiave, il valore e un campo che indica se la cella è vuota.  
Le due tabelle separate consentono il funzionamento della strategia Cuckoo, dove ogni elemento può essere spostato da una tabella all'altra in caso di collisione.

\begin{lstlisting}[language=C++, caption=Cuckoo Data Structure]
std::vector<Slot> table_cuckoo_1;
std::vector<Slot> table_cuckoo_2;
\end{lstlisting}

Durante l'inserimento, se la posizione calcolata tramite la prima funzione di hash è occupata, l'elemento presente viene espulso e 
reinserito nella posizione corrispondente della seconda tabella, garantendo che ogni elemento abbia sempre due possibili posizioni.


\section{Operazioni}

Le seguenti funzioni rappresentano le principali operazioni disponibili nella classe \texttt{HashTable}:

\begin{lstlisting}[language=C++, caption=Main Operations]
public:
    bool insert(const Tkey& key, const Tvalue& value);
    bool update(const Tkey& key, const Tvalue& value);
    bool remove(const Tkey& key);
    bool exist(const Tkey& key);
    std::optional<Tvalue> get(const Tkey& key);

private:
    bool resizeTable();
    void dynamicCopy();
    std::optional<std::size_t> find(const Tkey& key);
    std::optional<std::size_t> _insert(const Tkey& key, const Tvalue& value);
\end{lstlisting}

Le operazioni principali offerte dalla struttura dati sono:

\begin{itemize}
    \item \textbf{insert}: inserisce una nuova coppia chiave-valore nella tabella. Se la chiave è già presente, 
    l'operazione si comporta come una \texttt{update}.
    
    \item \textbf{update}: aggiorna il valore associato a una chiave esistente.
    
    \item \textbf{remove}: elimina un elemento dalla tabella, gestendo correttamente la struttura dati a seconda della strategia di collisione utilizzata.
    
    \item \textbf{exist}: verifica la presenza di una chiave all'interno della tabella.
    
    \item \textbf{get}: restituisce il valore associato a una chiave, se presente, utilizzando \texttt{std::optional} 
    per gestire il caso in cui la chiave non esista.
\end{itemize}

Tutte le operazioni di accesso si appoggiano alla funzione interna \texttt{find}, che implementa la logica di ricerca in base alla strategia di risoluzione 
delle collisioni selezionata (open addressing, chaining o cuckoo hashing).  
Questa funzione restituisce la posizione dell'elemento, se presente, oppure un valore nullo in caso contrario.

La funzione privata \texttt{\_insert} gestisce nel dettaglio la logica di inserimento, includendo la gestione delle collisioni e il posizionamento corretto 
dell’elemento nella struttura dati.  
Questa funzione viene utilizzata internamente da \texttt{insert} e può essere adattata a seconda della strategia adottata.

Il fattore di carico della tabella viene monitorato durante le operazioni di inserimento: se il valore supera la soglia \texttt{max\_load\_factor}, 
viene invocata la funzione \texttt{resizeTable}, che provvede ad aumentare la dimensione della tabella e a riorganizzare gli elementi.

Nel caso della strategia \texttt{DYNAMIC}, il ridimensionamento non avviene in modo immediato: la funzione \texttt{dynamicCopy} viene utilizzata per 
spostare progressivamente gli elementi dalla vecchia tabella a quella nuova, distribuendo il costo del resize su più operazioni e 
migliorando le prestazioni complessive.
